\section{Theoretical Analysis}
\label{sec:theory}

We chain five objects: constraint identifiability $\rightarrow$ projection perturbation $\rightarrow$ dynamic regret $\rightarrow$ predictability-conditioned advantage $\rightarrow$ failure-safe recovery.
Assumptions are collected in Appendix~\ref{app:impl} (cf.\ Assumps.~1--4 in the supplement markdown freeze).

\paragraph{Notation freeze.}
$\delta_t:=d_H(\hat{\mathcal{B}}_{t+1},\mathcal{B}_{t+1})$ (feasible-set mismatch);
$\chi_t:=d_H(\mathcal{B}_{t+1},\mathcal{B}_t)$ (constraint drift);
$\epsilon_t:=\|s_{t+1}-\hat s_{t+1}\|$ (state prediction error).
Capacity residual $\lvert\hat c_{t+1}-c_{t+1}\rvert$ upper-bounds $\delta_t$ for scalar budgets (Theorem~\ref{thm:1}) and is \emph{not} denoted $\delta_t$.

\subsection{Constraint Evolution Identifiability}
\begin{theorem}[Constraint evolution stability]
\label{thm:1}
For scalar budgets $\mathcal{B}(c)=\{x\ge 0:\mathbf{1}^\top x\le c\}$,
\begin{equation}
\delta_t
\;=\;
d_H\bigl(\mathcal{B}(c_{t+1}),\mathcal{B}(\hat c_{t+1})\bigr)
\;\le\;
\lvert c_{t+1}-\hat c_{t+1}\rvert.
\end{equation}
If $c=T(s)$ is $L_g$-Lipschitz and $\hat c_{t+1}=T(\hat s_{t+1})$, then
$\delta_t\le L_g\epsilon_t$.
\end{theorem}
Theorem~\ref{thm:1} establishes a controllable relationship between capacity / state residuals and feasible-set mismatch $\delta_t$, not that the predictor recovers the true future constraint.
Proof: Appendix~\ref{app:thm1}.

\begin{lemma}[Surrogate consistency]
\label{lem:surrogate}
Let $L_c=\lvert\hat c-c\rvert^2$ be the constraint regression loss and assume the scalar-budget geometry of Theorem~\ref{thm:1}.
Then
\begin{equation}
\mathbb{E}[\delta_t]
\;\le\;
\mathbb{E}\lvert\hat c_t-c_t\rvert
\;\le\;
\sqrt{\mathbb{E}[L_c]}.
\end{equation}
Hence minimizing the constraint prediction loss contracts the theoretical set-mismatch term appearing in Theorems~\ref{thm:1}--\ref{thm:2}.
This does \emph{not} claim that minimizing $L_c$ minimizes $\mathrm{Reg}_T$.
\end{lemma}
Proof: Appendix~\ref{app:surrogate}.

\subsection{Projection Error Propagation}
\begin{lemma}[Projection perturbation]
\label{lem:2}
Under bounded projection sensitivity, for any $y$,
\begin{equation}
\bigl\|\Pi_{\hat{\mathcal{B}}}(y)-\Pi_{\mathcal{B}}(y)\bigr\|
\;\le\;
C_\Pi\, d_H(\hat{\mathcal{B}},\mathcal{B}).
\end{equation}
Consequently, with $x_{t+1}=\Pi_{\hat{\mathcal{B}}_{t+1}}(y_t)$ and $x_{t+1}^\circ=\Pi_{\mathcal{B}_{t+1}}(y_t)$,
$\|x_{t+1}-x_{t+1}^\circ\|\le C_\Pi\delta_t$.
Under $\mu$-strong convexity,
$\|x_{t+1}^\star-x_t^\star\|\le\mu^{-1}(\omega_t+L_g\chi_t)$.
\end{lemma}
\begin{remark}
\label{rem:Cpi}
The constant $C_\Pi$ characterizes the sensitivity of the projection operator rather than a universal property of arbitrary convex sets.
This assumption is not required for arbitrary convex sets, but holds for the polyhedral feasible regions considered in this work; boundedness then follows from standard Hoffman-type error bounds, and often $C_\Pi=O(1)$.
\end{remark}
Proof: Appendix~\ref{app:lem2}.

\subsection{Dynamic Variational Regret}
Define $\mathrm{Reg}_T=\sum_{t=0}^{T-1}\bigl(F_t(x_t)-F_t(x_t^\star)\bigr)$ and
$P_T=\sum_{t}\|x_{t+1}^\star-x_t^\star\|$.
\begin{theorem}[Dynamic regret upper bound]
\label{thm:2}
SECDO admits a dynamic regret upper bound governed by objective drift, constraint evolution, and feasible-set mismatch:
\begin{equation}
\mathrm{Reg}_T
\;\le\;
O\Bigl(\sqrt{T(1+P_T)}\Bigr)
\;+\;
O\Bigl(\sum_t(\epsilon_t+\delta_t)\Bigr),
\end{equation}
with
\begin{equation}
P_T
\;\le\;
\frac1\mu\sum_t\bigl(\omega_t+L_g\chi_t\bigr).
\end{equation}
\end{theorem}
We do \emph{not} claim that SECDO achieves optimal regret.
Proof: Appendix~\ref{app:thm2}.

\subsection{Predictability-Conditioned Feasibility Improvement}
Let $V_R(t)\le\chi_t$ and $V_A(t)\le\delta_t$ be reactive / anticipatory future-violation certificates, and $\mathrm{PI}_t=\delta_t/(\chi_t+\varepsilon)$.
\begin{theorem}[Predictability-conditioned feasibility improvement]
\label{thm:3}
If $\mathrm{PI}_t<1$ (i.e., $\delta_t<\chi_t$), anticipatory projection provides a tighter \emph{constraint-violation certificate} when set mismatch is smaller than constraint drift:
\begin{equation}
\delta_t<\chi_t
\quad\Rightarrow\quad
\text{anticipatory feasibility certificate strictly tighter than reactive}.
\end{equation}
\end{theorem}
This is a statement about feasibility certificates, not about objective superiority, and not a claim that SECDO is always better.
Proof: Appendix~\ref{app:thm3}.

\subsection{Failure-Safe Recovery}
\begin{corollary}[Graceful degradation]
\label{cor:4}
With $\alpha_t=1/(1+\mathrm{PI}_t^2)$ and $c^{\mathrm{mix}}_t=\alpha_t\hat c_{t+1}+(1-\alpha_t)c_t$,
when $\mathrm{PI}_t\gg 1$ one has $\alpha_t\to 0$ and SECDO gracefully degrades toward reactive behavior.
Writing $I_t^{\mathrm{fail}}=\mathbf{1}\{\mathrm{PI}_t>1\}$,
\begin{equation}
\mathrm{Reg}_{\mathrm{SECDO}}
\;\le\;
\mathrm{Reg}_{\mathrm{Reactive}}
\;+\;
O\Bigl(\sum_t I_t^{\mathrm{fail}}\delta_t\Bigr).
\end{equation}
\end{corollary}
We do not claim full recovery to a failure-free trajectory.
Proof: Appendix~\ref{app:cor4}.
